\subsection{K-Waay and BatchReceive}

\kwaay is a post-quantum, X3DH-like deniable authenticated key-exchange
protocol constructed from KEMs, signatures, and a split-KEM
\cite{collins2024kwaay}. Its complete construction contains substantially
more cryptographic structure than is required for the question studied here.
We therefore focus on the receiver-side operation that processes several
sender-associated inputs within a single protocol action: \batchreceive.

In the protocol interface, a sender invokes \textsf{Send} using its local
state and a receiver prekey bundle, obtaining a session key and a protocol
message. When receiver ephemeral key material is reused, several sender
messages may correspond to the same receiver ephemeral state. The receiver
then invokes \batchreceive for that state with a vector of inputs
\cite[Sec.~4.1]{collins2024kwaay}. Writing the vector as

\begin{equation}
  S = \bigl((\mathit{pk}_{j_1},\mathit{prek}_{j_1},m_{j_1}),\ldots,
             (\mathit{pk}_{j_d},\mathit{prek}_{j_d},m_{j_d})\bigr),
  \qquad d \geq 1,
  \label{eq:kwaay-batch-input}
\end{equation}

each component contains the public key of the claimed sender, the claimed
sender prekey bundle, and the corresponding protocol message. For each
component, \batchreceive verifies the relevant prekey information and performs
the required decapsulation operations using the receiver's long-term and
ephemeral secret material. If a split-KEM decapsulation fails, the construction
rejects the batch as a whole by returning an all-$\bot$ result vector;
otherwise it returns the corresponding vector of derived-key results,
individual elements of which may still be $\bot$
\cite[Sec.~5.1, Fig.~10]{collins2024kwaay}.

The purpose of this description is not to reproduce the complete handshake or
its computational security proof. Rather, it identifies the protocol object
relevant to this work: one receiver-side invocation may process multiple
sender-associated components using a shared receiver-side state.

This grouping makes relations among components relevant to the execution.
Each position is associated with a claimed sender, a message, and a
receiver-side key result. Our analysis focuses on the boundary at which such
a collection of components is treated as one admissible batch. We use the
term \emph{batch admission} for this analytical boundary. It is terminology
introduced in this work and does not denote a separate algorithm in the
original \kwaay specification.

\subsection{The Stated Distinct-Party Condition}

The full version of the \kwaay construction explicitly states that, for a
given \batchreceive call, each element of $S$ is assumed to correspond to a
different party \cite[Sec.~5.1, after Fig.~10]{collins2024kwaay}. This is a
same-batch condition: it constrains the relationship among components
appearing within one \batchreceive invocation. It does not state that a party
participates only once globally, runs only one session, or produces only one
protocol message.

We refer to this condition as \distinctparty. This name is introduced by this
work; the original \kwaay paper states the condition in natural language. The
condition identifies a same-batch party relation but does not prescribe a
particular implementation mechanism for enforcing it. In particular, the
statement itself does not specify whether enforcement is performed by a
client, server, caller, or other batching component. Nor do we infer that a
deployed implementation necessarily omits such a check. Section~3 formally
defines the relation and the adversarial and evidentiary boundary under which
it is studied.

\subsection{Identity Coordinates}

To distinguish the party-level condition from other forms of uniqueness, we
represent the information relevant to batch admission using the abstract entry

\begin{equation}
  E_i=(A_i,\mathit{sid}_i,m_i),
  \label{eq:abstract-entry}
\end{equation}

where $A_i$ denotes a modeled protocol principal, $\mathit{sid}_i$ denotes the
sender-session/occurrence coordinate used in the analysis, and $m_i$ denotes
sender-produced protocol content. This tuple is an analytical abstraction
rather than the wire representation of a \kwaay input.

The notation $\mathit{sid}_i$ follows the corresponding coordinate in the
symbolic model and does not import session semantics beyond that encoding.
Likewise, $A_i$ should not be interpreted as an account record, public-key byte
string, prekey, or implementation object. It represents the protocol-principal
coordinate over which the distinct-party assumption is interpreted.

The analysis keeps the following coordinates separate.

\begin{table}[htbp]
  \centering
  \caption{Identity coordinates relevant to batch admission.}
  \label{tab:identity-coordinates}
  \small
  \begin{tabularx}{\linewidth}{@{}lXX@{}}
    \toprule
    Coordinate & What it distinguishes & What distinctness does not imply \\
    \midrule
    Party $A$ & A modeled protocol principal & A distinct message, occurrence,
      prekey, or key encoding \\
    Message $m$ & Sender-produced protocol content & A different originating party \\
    Occurrence $\mathit{sid}$ & A sender-session/occurrence used in the analysis
      & A different party \\
    Slot $i$ & A position within one batch & A different party or message \\
    Batch $B$ & One grouped receiver-side execution & Global uniqueness of its
      participants \\
    Receiver state $r$ & Receiver state shared by the grouped execution & The
      identity or number of contributing parties \\
    \bottomrule
  \end{tabularx}
\end{table}

These coordinates induce different equivalence relations. A single party may
participate through multiple protocol occurrences and produce distinct
messages. Consequently,

\begin{equation}
  A_i=A_j
  \quad\text{is consistent with}\quad
  \mathit{sid}_i\neq\mathit{sid}_j
  \ \land\ 
  m_i\neq m_j.
  \label{eq:same-party-different-message}
\end{equation}

Likewise, $i\neq j$ establishes only that two components occupy distinct batch
positions; it does not imply $A_i\neq A_j$. Batch identity and receiver state
establish a common processing context but do not determine the number of
distinct parties represented inside that context.

This distinction is also important when interpreting authentication evidence.
Evidence binding one protocol component to party $A$ can establish the origin
associated with that component. It does not, by itself, establish that $A$
occurs at most once among all components of the batch. Similarly, inequality
between two messages distinguishes protocol content but does not establish
inequality between their originating party coordinates.

\subsection{Motivating Identity Problem}

Consider the candidate two-position batch

\begin{equation}
  B^{\star}
  =\bigl((A,\mathit{sid}_1,m_1),
          (A,\mathit{sid}_2,m_2)\bigr),
  \qquad
  \mathit{sid}_1\neq\mathit{sid}_2,
  \quad m_1\neq m_2.
  \label{eq:motivating-composition}
\end{equation}

The two components occupy different positions, arise from different modeled
occurrences, and contain different messages, yet both are associated with the
same party coordinate $A$. A restriction that excludes only duplicate
messages therefore does not exclude this composition merely from
$m_1\neq m_2$. In contrast, the distinct-party condition excludes it because
both entries have the same party coordinate.

The example exposes the logical non-implication

\begin{equation}
  m_1\neq m_2
  \ \not\Longrightarrow\ 
  \mathsf{party}(E_1)\neq\mathsf{party}(E_2).
  \label{eq:message-does-not-imply-party}
\end{equation}

Composition~\eqref{eq:motivating-composition} is a motivating analytical
example rather than a machine-checked attack trace. In particular, it lies
outside executions satisfying \kwaay's stated distinct-party assumption.
Its purpose is to identify the semantic question studied in this work: what
changes when the same-batch relation over party identities is relaxed, and
whether a restriction defined over another coordinate can substitute for that
relation.

The threat model and formal variants used to answer these questions are
defined in Sections~3 and~4.

\subsection{Research Questions}

The preceding distinction motivates three research questions.

\paragraph{RQ1: invariant removal.}
Which same-batch party compositions become reachable when the
\distinctparty requirement is removed from an otherwise comparable admission
model?

\paragraph{RQ2: component-level consequence.}
Within the receiver-side abstraction used in the formal model, can relaxation
of the distinct-party condition result in multiple accepted components
associated, through the modeled origin relation, with one matching sender
occurrence?

Here, an accepted component denotes an event in our symbolic abstraction and
should not be confused with the single accept/reject status assigned to a
receiver session in the original \kwaay security model.

\paragraph{RQ3: coordinate substitution.}
Is a restriction defined only over message identity sufficient to imply the
same-batch party relation imposed by \distinctparty under the modeled
lifecycle, or can message-distinct components still share a party coordinate?

The explicit party-coordinate restriction used later is a positive control,
not a separate unknown: it provides a controlled comparison for the target
relation while the message-coordinate restriction tests the proposed
substitution.

These questions distinguish the primary party-level property from the
component-level behavior used to observe its consequences. They do not
presuppose a break of \kwaay's cryptographic primitives, a vulnerability in a
particular deployed implementation, or a unique mechanism by which the
distinct-party condition must be enforced. Section~3 defines the corresponding
attacker capabilities, assumptions, and non-goals; Section~4 describes their
realization in the symbolic model.
